% pdf-env-boxes.tex — עטיפת סביבות המשפט של הספר במסגרות צבעוניות ל-PDF, מותאם לצבעי HTML.
% משתמשים ב-mdframed ולא ב-tcolorbox: tcolorbox אינו ממקם נכון תחת כיוון RTL (התיבה נדחפת
% מחוץ לעמוד והטקסט גולש), ואילו mdframed עוטף נכון סביבות עבריות מימין לשמאל.
\usepackage{amsthm}
\usepackage{xcolor}
\usepackage[framemethod=default]{mdframed}

% ==== שמות בעברית ====
\providecommand{\definitionname}{הגדרה}
\providecommand{\examplename}{דוגמה}
\providecommand{\exercisename}{תרגיל}
\providecommand{\remarkname}{הערה}
\providecommand{\theoremname}{משפט}
\providecommand{\lemmaname}{למה}
\providecommand{\propositionname}{טענה}
\providecommand{\corollaryname}{מסקנה}
\renewcommand{\proofname}{הוכחה}

% ==== צבעים (בהתאם ל-CSS) ====
\definecolor{envDefinitionRGB}{RGB}{255,99,71}    % tomato
\definecolor{envExampleRGB}{RGB}{138,43,226}      % blueviolet
\definecolor{envExerciseRGB}{RGB}{30,144,255}     % dodgerblue
\definecolor{envRemarkRGB}{RGB}{90,90,90}         % gray
\definecolor{envTheoremRGB}{RGB}{255,165,0}       % orange
\definecolor{envProofRGB}{RGB}{60,179,113}        % mediumseagreen

% ==== סגנון בסיסי לכל המסגרות (mdframed) ====
% מסגרת צבעונית + רקע מוברש קל, שוברת עמודים, צמודה לרוחב הטקסט המלא.
\mdfdefinestyle{envbase}{%
  linewidth=0.6pt,
  leftmargin=0pt, rightmargin=0pt,
  innerleftmargin=8pt, innerrightmargin=8pt,
  innertopmargin=6pt, innerbottommargin=6pt,
  skipabove=10pt, skipbelow=10pt%
}

\makeatletter
% parent counter: ב-book לפי chapter, אחרת לפי section
\@ifundefined{c@chapter}{\def\qnumwithin{section}}{\def\qnumwithin{chapter}}

\newcommand{\EnsureNewTheorem}[2]{%
  \@ifundefined{#1}{%
    \newtheorem{#1}{#2}[\qnumwithin]%
  }{}%
}
\makeatother

\AtBeginDocument{%
  \theoremstyle{definition}
  \EnsureNewTheorem{definition}{\definitionname}
  \EnsureNewTheorem{example}{\examplename}
  \EnsureNewTheorem{exercise}{\exercisename}

  \theoremstyle{remark}
  \EnsureNewTheorem{remark}{\remarkname}

  \theoremstyle{plain}
  \EnsureNewTheorem{theorem}{\theoremname}
  \EnsureNewTheorem{lemma}{\lemmaname}
  \EnsureNewTheorem{proposition}{\propositionname}
  \EnsureNewTheorem{corollary}{\corollaryname}

  % ==== עטיפה במסגרות צבעוניות ====
  \surroundwithmdframed[style=envbase,linecolor=envDefinitionRGB,backgroundcolor=envDefinitionRGB!8]{definition}
  \surroundwithmdframed[style=envbase,linecolor=envExampleRGB,backgroundcolor=envExampleRGB!8]{example}
  \surroundwithmdframed[style=envbase,linecolor=envExerciseRGB,backgroundcolor=envExerciseRGB!8]{exercise}
  \surroundwithmdframed[style=envbase,linecolor=envRemarkRGB,backgroundcolor=envRemarkRGB!6]{remark}
  \surroundwithmdframed[style=envbase,linecolor=envTheoremRGB,backgroundcolor=envTheoremRGB!8]{theorem}
  \surroundwithmdframed[style=envbase,linecolor=envTheoremRGB,backgroundcolor=envTheoremRGB!8]{lemma}
  \surroundwithmdframed[style=envbase,linecolor=envTheoremRGB,backgroundcolor=envTheoremRGB!8]{proposition}
  \surroundwithmdframed[style=envbase,linecolor=envTheoremRGB,backgroundcolor=envTheoremRGB!8]{corollary}
  \surroundwithmdframed[style=envbase,linecolor=envProofRGB,backgroundcolor=envProofRGB!8]{proof}
}
